\section{A structural classification of positive polynomial experiments}
\label{sec:v18-structure}
We first treat finite detectors without an ordered latent space or a
prescribed resolution flag. Polynomial dependence below is on the commands,
not on the latent variable. In particular, the class includes every finite
detector with affine retention commands and arbitrary continuous cells.

\subsection{Experiments, queries, and the augmented moment matrix}
Let $X$ be a compact metric space and let $\mu$ be a probability on $X$.
Fix $2\le N<\infty$. At time $i$, a command belongs to
$U_i=[-1,1]^{b_i}$, where $b_i\ge1$, and the report belongs to a finite set $\mathcal Y_i$.
The likelihoods have the form
\begin{equation}\label{eq:v18-polynomial-detector}
 L_{i,y}(u,x)=\sum_{\alpha\in F_{i,y}}u^\alpha c_{i,y,\alpha}(x),
 \qquad c_{i,y,\alpha}\in C(X),\qquad
 L_{i,y}\ge\kappa>0,\qquad \sum_yL_{i,y}=1.
\end{equation}
Each $F_{i,y}$ is finite. The commands in the exploration experiment are
independent uniform random variables, independent of the latent variable
and of coding randomness. A command is read once. The persistent resource
is one of $M$ labels; the clock and the known experiment are read-only.
The meaning of this resource is exactly that of
Definition~\ref{def:finite-state}. In particular, a command-generation
seed that remembers the prefix is not free side information.

For $0\le n<N$, let $\mathscr V_n$ be the span in $C(X)$ of the
likelihoods of all specified report words in the remaining times
$n+1,\ldots,N$, with arbitrary fixed future commands. It is finite
dimensional, by \eqref{eq:v18-polynomial-detector}, and contains $1$:
summing over all future report words at any fixed future commands gives
$1$. Choose a finite list of such physical queries
$H_{n,1},\ldots,H_{n,q_n}$ spanning $\mathscr V_n$, and fixed weights
$w_{n,j}>0$ with sum one. A query is chosen independently after the
checkpoint label is formed. Only that query and its remaining trials
are executed. Its binary outcome is the indicator of its specified
report word. Thus squared-loss excess is the squared Euclidean error in
\[
 p_n(h)=\left(\sqrt{w_{n,j}}\,
             \frac{\mu(P_hH_{n,j})}{\mu P_h}\right)_{j=1}^{q_n}.
\]
Open-loop word queries suffice for this span; an adaptive future tree
only forms finite sums of the same products.

For a fixed past report word $a=(a_1,\ldots,a_n)$ set
\[
 P_a(\boldsymbol u,x)=\prod_{i=1}^nL_{i,a_i}(u_i,x),\qquad
 Z_a=\mu P_a,\qquad D_n=\sum_{i=1}^nb_i,
\]
and form the polynomial vector and its augmented matrix
\begin{equation}\label{eq:v18-augmented-matrix}
 \begin{split}
 Y_a(\boldsymbol u)&=
 \left(Z_a,\bigl(\sqrt{w_{n,j}}\,\mu(P_aH_{n,j})\bigr)_{j=1}^{q_n}\right)^t,\\
 \mathcal A_a(\boldsymbol u)&=
 [\,Y_a\mid\partial_1Y_a\mid\cdots\mid\partial_{D_n}Y_a\,],\\
 d_n(\mu)&=\max_{a\in\prod_{i\le n}\mathcal Y_i}
       \rank_{\R(\boldsymbol u)}\mathcal A_a-1.
 \end{split}
\end{equation}
The rank is over the rational-function field with the actual prior
moments as real coefficients. No genericity of the prior is implicit.
The initial state is a singleton; set $d_0=0$.

\begin{lemma}[Augmentation and normalization]\label{lem:v18-normalized-rank}
Let $Y=(Z,Zp)^t$ be a differentiable vector with $Z>0$. Then
\begin{equation}\label{eq:v18-normalized-rank}
 \rank[\,Y\mid DY\,]=1+\rank Dp.
\end{equation}
Consequently $d_n(\mu)$ is the maximum rank of the actual checkpoint
prediction map over interior command tuples and report words. It is
independent of the chosen spanning query coordinates, provided their
weights are positive.
\end{lemma}
\begin{proof}
Subtract $p_j$ times the first row from row $j+1$. The first column
becomes $(Z,0)^t$, and each derivative column becomes
$(\partial_iZ,Z\partial_ip)^t$. Subtracting
$(\partial_iZ/Z)$ times the first column removes its top entry.
These invertible operations prove \eqref{eq:v18-normalized-rank}.
A polynomial minor is either identically zero or nonzero on a dense
open subset of the interior command cube. Thus rational-function rank
is the maximum attained rank there. A different spanning query list
gives linear maps in both directions on the vector augmented by its
constant coordinate. After differentiation these maps identify the
ranks. Notice that no assertion $P\in\operatorname{im}DP$ was needed:
the added evidence column is removed by normalization, not treated as
an additional acquired direction.
\end{proof}

Write $R^{\mathrm{av}}_{M,n}$ for the optimal unconditional checkpoint
excess when an arbitrary encoder of the complete prefix may retain
$M$ labels, and $R^{\mathrm{max}}_{M,n}$ for its worst-history analogue.
Write $\mathfrak R^{\mathrm{av}}_{M,N}$ and
$\mathfrak R^{\mathrm{max}}_{M,N}$ for the infima of the corresponding
maximum over checkpoints when \emph{one} causal filter is used.
Independent encoder and decoder randomness is allowed in all infima.

\begin{theorem}[Product-rank classification]\label{thm:v18-rank-classification}
In every experiment \eqref{eq:v18-polynomial-detector}, for every fixed
prior $\mu$, the whole reachable prediction set at checkpoint $n$ is
compact semialgebraic of dimension $d_n(\mu)$. If $d_n(\mu)>0$, then
for all integers $M\ge1$
\begin{equation}\label{eq:v18-checkpoint-rank-law}
 c_nM^{-2/d_n}\le R^{\mathrm{av}}_{M,n}
 \le R^{\mathrm{max}}_{M,n}\le C_nM^{-2/d_n}.
\end{equation}
If $d_n=0$, the reachable set is finite, with at most
$\prod_{i\le n}|\mathcal Y_i|$ points; its checkpoint risk is zero once
that many labels are available.

Set $d_* =\max_{1\le n<N}d_n$. If $d_*>0$, then
\begin{equation}\label{eq:v18-causal-rank-law}
 cM^{-2/d_*}\le\mathfrak R^{\mathrm{av}}_{M,N}
 \le\mathfrak R^{\mathrm{max}}_{M,N}\le CM^{-2/d_*}.
\end{equation}
The upper bound is attained by one deterministic stage-compatible
filter with reachable representatives. If $d_*=0$, one finite filter is
exact at every checkpoint. These conclusions require no full-support
assumption on $\mu$, no Chebyshev system, and no nondegeneracy
hypothesis besides the rank computed in \eqref{eq:v18-augmented-matrix}.
For fixed detector and query lists, all positive-rank comparison
constants can be chosen uniformly on a compact set of prior moment
vectors on which the relevant ranks are constant.
\end{theorem}

\subsection{Whole images and actual history mass}
\begin{lemma}[Image dimension and a minorized chart]\label{lem:v18-image-mass}
For a fixed checkpoint, let $S_n$ be the union of all reachable vectors
$p_n(h)$. Its semialgebraic format is bounded in terms of the finite
polynomial data and horizon. Its dimension is $d_n$. If $d_n>0$, some
coordinate projection $\pi:\R^{q_n}\to\R^{d_n}$ and some cube $Q$
of positive side length satisfy
\begin{equation}\label{eq:v18-image-minorization}
 \pi_*\operatorname{Law}(p_n(h))\ \ge\ \beta\operatorname{Unif}(Q),
 \qquad \beta>0,
\end{equation}
under the stated exploration law, with reports included rather than
conditioned away.
\end{lemma}
\begin{proof}
For each word, the graph is specified by its command cube and the
polynomial equations $Z_ap_j=Y_{a,j+1}$. Here
$Z_a\ge\kappa^n$ on that cube. Projection gives a semialgebraic set of
bounded format. Its compactness follows from the continuous rational
map on a compact cube. The finite union over words has the same
properties. Prior integrals are simply real coefficients in these
equations; neither $X$ nor the prior has been assumed semialgebraic.

For clarity, the real-geometric dimension fact used here is the
following standard consequence of semialgebraic constant-rank
stratification \cite{BochnakCosteRoy}: a rational map with nonvanishing
denominator on a cube has image dimension equal to its largest
interior differential rank. Indeed, stratify its domain, including the
boundary faces, into finitely many smooth semialgebraic pieces on which
the restricted map has constant rank. The constant-rank theorem bounds
each image dimension by that rank. Every restricted rank is at most
the maximum ambient rank, which is also the maximum interior rank:
a nonzero minor at a boundary point stays nonzero at nearby interior
points. Conversely, a point of maximum interior rank has a neighborhood
whose image contains a submanifold of that dimension. Applying this
fact and Lemma~\ref{lem:v18-normalized-rank} proves $\dim S_n=d_n$.

Choose a word and an interior tuple $u_*$ where the rank is $d_n$, and
choose $d_n$ output rows with independent differentials. Put $z=\pi p_n$
and $J=Dz(u_*)$. Let the rows of $K$ be an orthonormal basis of
$\ker J$. The square map
\[
 F(u)=(z(u),K(u-u_*))
\]
has an invertible derivative at $u_*$. The inverse-function theorem
provides a diffeomorphism over a product of a $d_n$-cube and a
$(D_n-d_n)$-cube. Shrinking these cubes makes the forward Jacobian
bounded, say by $B$ in operator norm. Thus its inverse determinant is
at least $B^{-D_n}$. In command coordinates the subprobability density
of this very report word is $Z_a/\operatorname{vol}(\prod U_i)$,
bounded below by $\kappa^n/\operatorname{vol}(\prod U_i)$.
Changing variables and integrating over the entire kernel-coordinate
cube gives a positive lower density on the first cube. This proves
\eqref{eq:v18-image-minorization}. Restricting to a local section instead
would discard its ambient command mass; we have not done so.
\end{proof}

\begin{proof}[Proof of Theorem~\ref{thm:v18-rank-classification}]
First let $d=d_n>0$. The vector $p_n$ lies in a fixed unit box.
Lemma~\ref{lem:tame-rectangle}, applied to the whole set of
Lemma~\ref{lem:v18-image-mass}, gives
$\mathcal N(S_n,\varepsilon)\le A\varepsilon^{-d}$ for
$0<\varepsilon\le1$. The constant depends only on its bounded format
and ambient dimension. Replacing every nonempty covering center by a
point of $S_n$ at most doubles its radius. Consequently $M$ reachable
representatives give error at most $C_nM^{-1/d}$; for the finitely many
small budgets in this deduction increase $C_n$ and use one point.
Encoding a nearest representative proves the checkpoint upper bound.

For the lower bound, project the $M$ arbitrary prediction centers to
the coordinates in \eqref{eq:v18-image-minorization}. If the cube has
volume $v$ and $\omega_d$ is the volume of the unit $d$-ball, the union
of their radius-$r$ balls occupies at most $M\omega_dr^d$ of the cube.
Take $r=(v/(2M\omega_d))^{1/d}$. At least half the uniform mass remains
at distance at least $r$, so the unconditional squared error is at
least $\beta r^2/2$. Conditional averaging first makes each randomized
decoder a deterministic vector for each label. A randomized assignment
to these vectors cannot improve their pointwise minimum distance.
An independent public coding seed can be fixed and then averaged;
its conditioning does not change the exploration law. This proves the
lower bound with exactly the permitted randomization. Worst-history
risk is at least unconditional risk.

When $d_n=0$, every fixed-word prediction map has derivative zero on
the connected interior cube, and is therefore constant there and on
its closure. There is at most one point per word. This conclusion
allows several distinct points; zero differential rank does not assert
that one label suffices in a general finite-report experiment.

We next construct one causal filter. For fixed $u,y$, multiplication
by $L_{n+1,y}(u,\cdot)$ maps $\mathscr V_{n+1}$ into $\mathscr V_n$.
Since $1\in\mathscr V_{n+1}$, it also puts the next likelihood in
$\mathscr V_n$. Choose expansions in the fixed spanning list at time
$n$. Their coefficients are polynomial in $u$ and bounded on $U_{n+1}$:
this follows by choosing a fixed linear right inverse to that spanning
list on the finite-dimensional coefficient space. Hence the next
prediction vector is a well-defined rational function
\begin{equation}\label{eq:v18-update}
 T_n(p,u,y)_j=\sqrt{w_{n+1,j}}\,
 \frac{\nu_p(L_{n+1,y}(u,\cdot)H_{n+1,j})}
      {\nu_p L_{n+1,y}(u,\cdot)}.
\end{equation}
Here $\nu_p$ denotes any posterior realizing $p$; its action on
$\mathscr V_n$ is determined by $p$ and $\nu_p1=1$, so the expression
is independent of the realization.

The denominator in \eqref{eq:v18-update} is at least $\kappa$, not only
on $S_n$ but also on line segments between its points: such a segment
is represented on $\mathscr V_n$ by the corresponding mixture of
posteriors. Bounded coefficients and the quotient rule therefore give
\[
 \|T_n(p,u,y)-T_n(p',u,y)\|\le L_n\|p-p'\|
 \quad(p,p'\in S_n),
\]
with $L_n<\infty$ independent of histories and commands. If $p$ is
reachable, so is $T_n(p,u,y)$, since the appended report has strictly
positive probability.

For $M$ large enough, choose at stage $n$ a reachable quantizer $Q_n$
of radius at most $A_nM^{-1/d_n}$ when $d_n>0$, and store every point
when $d_n=0$. With $d_* >0$ all these radii are at most
$AM^{-1/d_*}$. Store only the representative index and use
\[
 \widehat p_{n+1}=Q_{n+1}(T_n(\widehat p_n,u_{n+1},y_{n+1})).
\]
This is an admitted transition using only the preceding index and the
current input. Starting with the known exact initial vector, its errors
satisfy
\[
 e_0=0,\qquad e_{n+1}\le L_ne_n+AM^{-1/d_*}.
\]
Finite iteration gives $e_n\le B_NM^{-1/d_*}$ at every checkpoint.
For the finitely many excluded budgets, the squared prediction loss is
at most one, so increasing the constant supplies the same assertion.
The decoder uses the actual probabilities of the representative. This
proves the simultaneous upper bound without retaining or reconstructing
an exact discarded prefix. Every causal filter is in particular an
$M$-message prefix encoder at a checkpoint attaining $d_*$; the lower
bound already proved there gives \eqref{eq:v18-causal-rank-law}.
If every $d_n=0$, retain all finitely many reachable states at each
stage instead; the same exact updates give a finite exact filter.

Finally fix the detector and queries and vary their finite prior moment
vector in a compact constant-rank set. All first and second command
derivatives of the normalized maps are uniformly bounded, because
evidence is bounded below by $\kappa^n$. At each parameter choose a
nonzero $d_n$-minor and an interior tuple. Continuity preserves a
positive minor on a parameter neighborhood. A finite subcover gives
uniform chart radii and lower mass by the preceding inverse-function
argument. The whole-image format bound was already coefficient
independent. The update coefficients are fixed by the detector and
queries, and the denominator is still at least $\kappa$. These
observations give all asserted uniform constants.
\end{proof}

The theorem identifies the exponent directly from the experiment, not
from a postulated scale representation of its risk. It makes no claim
that the integer $d_n$ determines a uniform profile when ranks change.
The covariance classification in Section~\ref{sec:v18-affine}, and the
collision and contrast classifications later, supply that additional
information in their respective classes. The results also distinguish
semialgebraic image dimension from the dimension of a Euclidean space
into which the entire exact state set might embed.

\subsection{Algebraic rank strata in the moment body}
Choose a basis $1,g_1,\ldots,g_s$ of the finite-dimensional space
generated by all coefficient functions needed to form all vectors
$Y_a$ at the fixed horizon. Write
\[
 m(\mu)=(\mu g_1,\ldots,\mu g_s),\qquad
 K=\operatorname{conv}\{(g_1(x),\ldots,g_s(x)):x\in X\}.
\]
This convex hull is compact. Linear independence with $1$ means that
$K$ has nonempty interior in $\R^s$; the case $s=0$ is read in its
zero-dimensional affine space. Entries of every $\mathcal A_a$ are
polynomials in commands whose coefficients are affine functions of $m$.

\begin{theorem}[Exact moment strata and generic maximal acquisition]
\label{thm:v18-moment-strata}
For every checkpoint $n$ and integer $k\ge0$ there is an explicit finite
list of polynomials $\mathcal F_{n,k}$ in $m$ such that
\begin{equation}\label{eq:v18-stratum}
 d_n(\mu)\le k\quad\Longleftrightarrow\quad
 f(m(\mu))=0\ \hbox{for every }f\in\mathcal F_{n,k}.
\end{equation}
The list consists of the command coefficients of all $(k+2)$-minors
of the matrices \eqref{eq:v18-augmented-matrix}, over all report words.
Thus the complete checkpoint and causal memory exponents are determined
by finitely many polynomial equalities in the actual prior moments.

The moments of full-support probabilities are precisely $\operatorname{int}K$.
On this set there is a relatively open dense complement of a proper
algebraic set on which all checkpoint ranks simultaneously take their
maximum values. These generic ranks can be reached by arbitrarily small
bounded-density perturbations of any full-support prior. Lower-rank
loci, when present, are exactly the algebraic loci in
\eqref{eq:v18-stratum}, intersected with the moment body.
\end{theorem}
\begin{proof}
The equivalence is the polynomial rank criterion: a command polynomial
vanishes on the interior cube exactly when every one of its coefficients
vanishes. The normalization rank is one smaller by
Lemma~\ref{lem:v18-normalized-rank}. An empty list of minors has its
usual vacuous meaning when $k$ exceeds the possible rank. This gives a
finite algebraic test using exact moments, not an equality test from
inexact numerical moment names.

Suppose $\mu$ has full support and $m(\mu)$ were on a proper supporting
hyperplane of $K$. Its defining affine function would give a continuous
nonnegative function on $X$, with integral zero under $\mu$, but not
identically zero. Positivity on a nonempty open set contradicts full
support. Hence $m(\mu)\in\operatorname{int}K$. Conversely, take an
interior $m$ and any fixed full-support $\mu_0$. Such a probability
exists on a compact metric space, for example by assigning positive
summable masses to a countable dense set. For sufficiently small
$\epsilon>0$,
\[
 m'=\frac{m-\epsilon m(\mu_0)}{1-\epsilon}\in K.
\]
By the definition of $K$ and finite-dimensional convexity, $m'$ is a
finite convex combination of points $g(x)$ and hence is the moment of
a probability $\nu$. The probability
$\epsilon\mu_0+(1-\epsilon)\nu$ has full support and moment $m$.

For each checkpoint of positive maximal rank choose one nonzero
coefficient of a maximal nonzero minor. Its polynomial is not
identically zero on $\R^s$: $K$ has interior. The product of the
finitely many chosen polynomials is nonzero. Away from its zero set
all maximal ranks are attained simultaneously. Checkpoints whose
maximal rank is zero impose no condition. A nonzero real polynomial
has a zero set of empty interior and Lebesgue measure zero, giving the
claimed dense open set. No semialgebraicity of $K$ itself is asserted.

For the final assertion fix a full-support $\mu$ and put $m=m(\mu)$.
The covariance matrix $G=\mu[(g-m)(g-m)^t]$ is positive definite:
zero variance of a linear combination would make it constant on $X$,
contrary to the independence of $1,g_1,\ldots,g_s$. For sufficiently
small $t\in\R^s$ the probability
\[
 d\mu_t=(1+t^t(g-m))\,d\mu
\]
is positive, has total mass one, has full support, and has moment
$m+Gt$. These moments fill a neighborhood of $m$, while the densities
tend uniformly to one. Every such neighborhood meets the generic
maximal-rank set just constructed.
\end{proof}

Together, Theorems~\ref{thm:v18-rank-classification} and
\ref{thm:v18-moment-strata} give a necessity-and-sufficiency description
of the memory exponent throughout this class, including exceptional
priors. The algebraic test does not presuppose a Chebyshev minorization,
a rectangular risk curve, or a known attainable geometry. What it does
not encode is a sharp anisotropic profile near an exceptional locus.
We now obtain such a profile, uniformly through all rank losses, for a
whole class of affine acquisitions.
